\documentclass[12pt]{exam}
\usepackage[margin=0.1in]{geometry}

\usepackage{amsmath,amssymb}
\usepackage{graphicx}
\usepackage{enumitem}
\usepackage{multicol}
\usepackage{xcolor}
\usepackage{tikz}
\usetikzlibrary{arrows.meta,positioning,shapes.callouts}

\newcommand{\tfq}[3]{%
  \question[#1]\textbf{T/F:} #2%
  % Only show the choices when printing answers/keys
  \ifprintanswers
    \par\begin{oneparchoices}
      #3
    \end{oneparchoices}
  \fi
}



\makeatletter
\newcommand{\TFQ@choices}{} % temp holder

\newenvironment{tfqverb}[2]{%
  \def\TFQ@choices{#2}% store choices for use at end
  \question[#1]\textbf{T/F:}%
}{%
  \ifprintanswers
    \par\begin{oneparchoices}\TFQ@choices\end{oneparchoices}
  \fi
}
\makeatother

\printanswers
% optional styling for the correct choice:
\CorrectChoiceEmphasis{\bfseries}    % make the correct choice bold
% or e.g. \CorrectChoiceEmphasis{\color{red}\bfseries}

\usepackage{fancyvrb}
\newsavebox{\verbbox}

\begin{document}


\begin{center}
\textbf{\Large Test 2 — Database Management Systems}

\textbf{CSC 3300 \quad Spring 2026}
\end{center}

\begin{questions}
%---------------------------------TRUE/FALSE---------------------------------
\section*{Part I — True/False}
\tfq{3}{If the attribute \texttt{room\_number} of the \texttt{classroom} table was removed from the \texttt{PRIMARY KEY} (but not removed from the table), then in each building on campus there could be at most one classroom.}{\CorrectChoice True \choice False}
\tfq{3}{A tuple is inserted into the \texttt{instructor} relation only if the value of the \textit{ID} attribute of this tuple is equal to the value of the \textit{i\_id} attribute of some tuple in the \texttt{advisor} relation.}{\choice True \CorrectChoice False}
\begin{tfqverb}{3}{\choice True \CorrectChoice False}
Below query returns at most one tuple on any instance of the University database.

\begin{verbatim}
select course_id
from course
where course_id = 128;
\end{verbatim}
\end{tfqverb}
\tfq{2}{The \texttt{section} relation has exactly 3 foreign keys.}{\choice True \CorrectChoice False}

\begin{tfqverb}{3}{\CorrectChoice True \choice False}
The following query returns IDs of students along with IDs of instructors who taught that student in more than two course sections.

\begin{verbatim}
select takes.ID, teaches.ID 
from takes join teaches using (course_id, sec_id, semester, year)
group by takes.ID, teaches.ID 
having count(teaches.ID) > 2
\end{verbatim}
\end{tfqverb}

%---------------------------------MC---------------------------------
\section*{Part II — Multiple choice (choose one)}

% \begin{minipage}{\linewidth}
% \question[3]
% What is the minimal number of tables that need to be reached out to in order to retrieve \textit{IDs} and \textit{names} of students and \textit{ID}'s and \textit{names} of their teachers (past or present)?
% \begin{choices}
%   \choice 3
%   \CorrectChoice 4
%   \choice 5
%   \choice 2
% \end{choices}
% \end{minipage}\vspace{1em}

\begin{minipage}{\linewidth}
\question[3]
What is the minimal number of tables that need to be reached out to in order to retrieve \textit{IDs} and names of instructors along with \textit{IDs} and titles of courses they taught, that come from a different department than the one they are associated with?

E.g. If an instructor with id 5 and name Green from Math Dept. taught a course with id CSC3300 and title DBMS from the CS Dept., then the below tuple should be included in the result of the query.

\begin{center}
\begin{tabular}{|c|c|c|c|}
\hline
\textit{ID} & \textit{name} & \textit{course\_id} & \textit{title} \\
\hline
5 & Green & CSC3300 & DBMS \\
\hline
\end{tabular}
\end{center}

\begin{checkboxes}
  \choice 4
  \choice 2
  \CorrectChoice 3
  \choice 5
\end{checkboxes}

\end{minipage}\vspace{1em}

%---------------------------------MS---------------------------------
\section*{Part III — Multiple select (select all that apply)}

\begin{minipage}{\linewidth}
\question[4] Choose superkeys of the relation \texttt{takes}.\par\vspace{0.5em}

\begin{checkboxes}
\choice $\{\textit{ID}, \textit{course\_id}, \textit{sec\_id}, \textit{semester}, \textit{year}\}$
\CorrectChoice $\{\textit{ID}, \textit{grade}\}$
\CorrectChoice $\{\textit{course\_id}, \textit{sec\_id}, \textit{semester}, \textit{year}\}$
\choice $\{\textit{ID}, \textit{course\_id}, \textit{sec\_id}, \textit{semester}, \textit{year}, \textit{grade}\}$
\end{checkboxes}
\end{minipage}\vspace{1em}


\begin{minipage}{\linewidth}
\question[4]
An ISBN (International Standard Book Number) is a unique identifier for books, allowing them to be easily identified and tracked in databases and systems worldwide. An ISBN is assigned to each edition and variation (e.g., hardcover, paperback, e-book) of a book.

Consider the following specification of user requirements for a Library Database: For each book copy, the following data is stored: title (\textit{title}), edition (\textit{edition}), number of pages (\textit{pages\_no}), barcode (\textit{barcode}), variation (e.g., hardcover, paperback) (\textit{variation}), and ISBN (\textit{isbn}). When a member checks out a book copy, its barcode is scanned and linked to the member’s account. When the book copy is returned, the barcode is scanned again, and the book is marked as available for borrowing.

To organize this data, assume the following relational database schema:

\begin{center}
\texttt{BOOK (title, edition, pages\_no, variation, isbn)}\\
\texttt{BOOK\_COPY (isbn, barcode)}
\end{center}

Assume that the attribute \textit{isbn} in the BOOK table holds the same data (has the same meaning) as the attribute \textit{isbn} in the BOOK\_COPY table.

Which of the following statements are true?\par\vspace{0.5em}

\begin{checkboxes}
\CorrectChoice $\{\textit{isbn}\}$ is a candidate key in the BOOK table.
\CorrectChoice If $\{\textit{isbn}\}$ was the primary key of BOOK, then $\{\textit{isbn}\}$ could be a foreign key in the BOOK\_COPY table.
\CorrectChoice $\{\textit{barcode}\}$ is a candidate key in the BOOK\_COPY table.
\CorrectChoice If multiple editions or variations of the same book exist, the title may appear in multiple tuples in the BOOK table.
\end{checkboxes}
\end{minipage}\vspace{1em}


\begin{minipage}{\linewidth}
\question[8]
IDs, names, and major departments (\texttt{dept\_name}) of students who have not received a grade of \texttt{A} in any of their courses.\par\vspace{0.5em}

\begin{checkboxes}
\CorrectChoice
\begin{verbatim}
select ID, name, dept_name
from student
where not exists (
  select *
  from takes
  where student.ID = takes.ID and grade='A')
\end{verbatim}

\CorrectChoice
\begin{verbatim}
select ID, name, dept_name
from student
where 'A' <> all (
  select grade
  from takes
  where student.ID = takes.ID)
\end{verbatim}

\choice
\begin{verbatim}
select ID, name, dept_name
from student
where exists (
  select *
  from takes
  where grade <> 'A')
\end{verbatim}

\choice
\begin{verbatim}
select ID, name, dept_name
from student
where not exists (
  select *
  from takes
  where student.ID = takes.ID and not (grade='A'))
\end{verbatim}

\end{checkboxes}
\end{minipage}

\begin{minipage}{\linewidth}
\question[8]
Which query or queries return \textit{IDs}, \textit{names}, and \textit{dept\_names} of students that didn’t take any course sections?\par\vspace{0.5em}

\begin{checkboxes}
\choice
\begin{verbatim}
select student.ID, name, dept_name
from student, takes
where student.ID <> takes.ID
\end{verbatim}

\CorrectChoice
\begin{verbatim}
select ID, name, dept_name
from student
where (select count(*) 
       from takes 
       where takes.ID = student.ID) = 0
\end{verbatim}

\CorrectChoice
\begin{verbatim}
select ID, name, dept_name
from student
where not exists (select * 
                  from takes 
                  where ID = student.ID)
\end{verbatim}

\CorrectChoice
\begin{verbatim}
select ID, name, dept_name
from student
where ID not in (select ID from takes)
\end{verbatim}

\end{checkboxes}
\end{minipage}\vspace{1em}

\begin{minipage}{\linewidth}
\question[8]
Which queries retrieve IDs and names of students who have taken at least one Comp. Sci. course.\par\vspace{0.5em}

\begin{checkboxes}
\choice
\begin{verbatim}
select ID, name
from takes natural join course natural join student
where dept_name = 'Comp. Sci.'
group by ID, name
having count(*) >= 1
\end{verbatim}

\CorrectChoice
\begin{verbatim}
select ID, name
from student as s
where (select count(*)
      from takes natural join course
      where dept_name = 'Comp. Sci.'
      and s.ID = takes.ID) >= 1
\end{verbatim}

\CorrectChoice
\begin{verbatim}
select student.ID, name
from (select ID
      from takes natural join course
      where dept_name = 'Comp. Sci.'
      group by ID
      having count(*) >= 1) as temp, student
where temp.ID = student.ID
\end{verbatim}

\CorrectChoice
\begin{verbatim}
select ID, name
from (select ID
      from takes natural join course
      where dept_name = 'Comp. Sci.'
      group by ID
      having count(*) >= 1) as temp natural join student
\end{verbatim}

\end{checkboxes}
\end{minipage}\vspace{1em}

\begin{minipage}{\linewidth}
\question[8]
Which of the following queries return the names of \textbf{all} buildings at the university? 
Ignore any repeated building names in the results.\par\vspace{0.5em}

\begin{checkboxes}
\choice
\begin{verbatim}
select building
from department natural join classroom;
\end{verbatim}

\CorrectChoice
\begin{verbatim}
select building
from department natural join classroom
union
select building
from department
where building not in (select building from classroom)
union
select building
from classroom
where building not in (select building from department);
\end{verbatim}

\choice
\begin{verbatim}
select department.building
from department join classroom;
\end{verbatim}

\CorrectChoice
\begin{verbatim}
select building
from department
left outer join classroom using (building)
union
select building
from department
right outer join classroom using (building);
\end{verbatim}

\end{checkboxes}
\end{minipage}\vspace{1em}


\end{questions}
\end{document}